\begin{trapbox}
	\begin{description}[style=nextline, leftmargin=2.8em, labelsep=0.5em, itemsep=0.35em]
		\item[\textbf{\textcolor{CTrap}{易错点一（公式混淆）}}]
		      错误地认为 $E\xi=np(1-p)$。
		\item[\textbf{\textcolor{CTrap}{正确理解（正确公式）：}}]
		      二项分布期望为 $E\xi=np$，方差为 $D\xi=np(1-p)$。
		\item[\textbf{\textcolor{CTrap}{错因分析（记忆模糊）：}}]
		      期望与方差公式混淆，或受方差形式影响。
	\end{description}

	\medskip
	\begin{description}[style=nextline, leftmargin=2.8em, labelsep=0.5em, itemsep=0.35em]
		\item[\textbf{\textcolor{CTrap}{易错点二（把“不是二项分布”误认为“一定不能求期望”）}}]
		      错误地认为只要不是独立重复试验，就一定不能得到 $E\xi=np$。

		\item[\textbf{\textcolor{CTrap}{正确理解（区分分布与期望）：}}]
		      若要直接判定 $\xi\sim B(n,p)$，必须满足独立重复试验，且每次成功概率均为 $p$。
		      但若只求成功次数的期望，独立性不是必要条件。设
		      \[
			      \xi=X_1+X_2+\cdots+X_n,
		      \]
		      其中 $X_i$ 表示第 $i$ 次试验是否成功，则
		      \[
			      E\xi=\sum_{i=1}^n EX_i=\sum_{i=1}^n p_i.
		      \]
		      如果各次边际成功概率都等于 $p$，则仍有 $E\xi=np$。

		\item[\textbf{\textcolor{CTrap}{典型提醒（不放回抽样）：}}]
		      不放回抽样的成功次数一般服从超几何分布，不是二项分布；
		      若总体中成功个数为 $M$，总体容量为 $N$，成功比例为 $p=\dfrac{M}{N}$，抽取 $n$ 次，则
		      \[
			      E\xi=n\frac{M}{N}=np.
		      \]

		\item[\textbf{\textcolor{CTrap}{错因分析（混淆公式来源）：}}]
		      二项分布公式 $E\xi=np$ 可以直接套用，但 $E\xi=np$ 这个形式并非二项分布独有。
		      求期望时更本质的方法是指示变量分解与期望线性性。
	\end{description}

	\medskip
	\begin{description}[style=nextline, leftmargin=2.8em, labelsep=0.5em, itemsep=0.35em]
		\item[\textbf{\textcolor{CTrap}{易错点三（平方期望误用）}}]
		      错误地认为 $E(\xi^2)=(E\xi)^2=(np)^2$。
		\item[\textbf{\textcolor{CTrap}{正确理解（正确公式）：}}]
		      利用方差公式，
		      \[
			      E(\xi^2)=D\xi+(E\xi)^2=np(1-p)+(np)^2.
		      \]
		\item[\textbf{\textcolor{CTrap}{错因分析（线性混淆）：}}]
		      将期望的线性性质误推广到平方运算，注意 $E(\xi^2)\neq(E\xi)^2$。
	\end{description}
\end{trapbox}